\section{Knowledge Integration}

Since this bachelor's thesis belongs to the software engineering specialization, the knowledge acquired in various specialization subjects is applied throughout the project lifecycle. Below is a list of each of these subjects, along with their specific contributions to this project

\vspace{0.2cm}
\textbf{Web Services and Applications (ASW)}
\vspace{-0.1cm}

Provided that the project is centered on a web application, this is undoubtedly the first subject to address. The key outcome of this subject is a solid understanding of RESTful API design. At the same time, hands-on experience with Ruby on Rails has strengthened the author's skills in database management and migrations. Finally, the introduction of the Swagger tool has also improved API testing and documentation. As a result, these concepts have been actively applied throughout the project.

\vspace{0.2cm}
\textbf{Software Engineering Projects (PES)}
\vspace{-0.1cm}

PES is also a key component of this project. This subject gave the author the first opportunity to develop a software application from the ground up, covering all necessary phases of the software lifecycle. Through this experience, the subject highlights that software engineering goes beyond coding, incorporating essential planning and management activities to ensure project success.

On top of that, it enabled the practical application of Agile methodologies, which fostered pragmatic habits such as continuous documentation and regular updates to the product backlog. Additionally, it provided the opportunity to build a decoupled frontend–backend application and work with unfamiliar frameworks, significantly strengthening the author's expertise in React. 

The introduction of new tools has also been beneficial for the author and has therefore enabled their application in this project. For example, editor extensions that support coding best practices and Git version control.

\vspace{0.2cm}
\textbf{Software Project Management (GPS)}
\vspace{-0.1cm}

GPS is the first subject that introduced the author to agile methodology concepts, one of the most relevant approaches in the current software market. It also taught how to create an initial timeline and budget estimate. Furthermore, it covered other common software project concepts, such as risk analysis and the Gantt diagram.

\vspace{0.2cm}
\textbf{Requirements Engineer (ER)}
\vspace{-0.1cm}

The project's context study, requirements analysis, and use case diagrams are all documented in accordance with the concepts introduced in this specific subject.

\vspace{0.2cm}
\textbf{Software Engineering Introduction (IES) and Software Architecture (AS)}
\vspace{-0.1cm}

These two subjects introduced the basics for building UML diagrams and designing layered architecture.

\vspace{0.2cm}
\textbf{Database (BD) and Database Design (DBD)}
\vspace{-0.1cm}

As the name suggests, these subjects helped the author design a secure and optimized database, write efficient queries, and become familiar with SQL databases.

\vspace{0.2cm}
\textbf{Cybersecurity Management (GCS)}
\vspace{-0.1cm}

The theory presented in this subject raised the author's awareness of the wide variety of cyberattack strategies and underscored the importance of preventing any foreseeable cybersecurity risks.

\vspace{0.2cm}
\textbf{Informatics' Social and Environmental Aspects (ASMI)}
\vspace{-0.1cm}

While the previous subjects provided more technical knowledge, this one focuses on non-technical aspects. Specifically, it has significantly enriched the project's sustainability analysis.
